% Part: sets-functions-relations
% Chapter: functions
% Section: kinds

\documentclass[../../../include/open-logic-section]{subfiles}

\begin{document}

\olfileid{sfr}{fun}{kin}
\olsection{फलनांचे प्रकार}

\begin{explain}
आपल्याला वारंवार भेटणाऱ्या फलनांच्या काही प्रकारांचे वर्गीकरण करून
घेणे उपयोगी ठरेल.

सुरुवातीला, सहप्रांतातील प्रत्येक सदस्य फलनाचे मूल्य असतो असा गुणधर्म
असलेल्या फलनांचा विचार करू. अशा फलनांना !!{surjective} म्हणतात.
त्यांची आकृती \olref{fig:surjective} प्रमाणे काढता येते.

\begin{figure}
  \olasset{assets/diagrams/surjective.tikz}
  \caption{!!^a{surjective} फलनात सहप्रांतातील प्रत्येक !!{element}
    हे फलनाचे मूल्य असते.}
  \ollabel{fig:surjective}
\end{figure}
\end{explain}

\begin{defn}[!!^{surjective} फलन]
$f \colon A \rightarrow B$ हे फलन \emph{!!{surjective}} असते तेव्हा
आणि केवळ तेव्हाच, जेव्हा $B$ हा~$f$ ची व्याप्तीही असतो. म्हणजे,
प्रत्येक $y \in B$ साठी~$f(x) = y$ असणारा किमान एक $x \in A$
असतो. चिन्हांत:
\[
  (\forall y \in B)(\exists x \in A)f(x) = y.
\]
अशा फलनाला $A$ पासून $B$ पर्यंतचे !!a{surjection} म्हणतो.
\end{defn}

\begin{explain}
$f$ हे !!a{surjection} आहे असे दाखवायचे असल्यास, $f$ च्या सहप्रांतातील
प्रत्येक वस्तू कोणत्या तरी आदान $x$ साठी $f(x)$ चे मूल्य आहे असे दाखवावे लागते.

कोणत्याही फलनापासून एक !!a{surjection} \emph{निर्माण होते}, हे लक्षात घ्या.
कारण $f \colon A \to B$ हे फलन दिले असल्यास, $f' \colon A \to \ran{f}$
ची व्याख्या $f'(x) = f(x)$ अशी करू. $\ran{f}$ ची \emph{व्याख्याच}
$\Setabs{f(x) \in B}{x \in A}$ अशी असल्यामुळे हे फलन $f'$ निश्चितच
!!a{surjection} असते.
\end{explain}

\begin{explain}
कोणतेही फलन प्रत्येक शक्य आदानाशी एकमेव प्रदान जोडते. पण वेगवेगळ्या
आदानांशी कधीच एकच प्रदान न जोडणारी फलनेही असतात. अशा फलनांना
!!{injective} म्हणतात. त्यांची आकृती \olref{fig:injective} प्रमाणे काढता येते.
\begin{figure}
  \olasset{assets/diagrams/injective.tikz}
  \caption{!!^a{injective} फलन कधीच दोन वेगवेगळ्या फलसाधकांशी
    एकच मूल्य जोडत नाही.}
  \ollabel{fig:injective}
\end{figure}
\end{explain}

\begin{defn}[!!^{injective} फलन]
$f \colon A \rightarrow B$ हे फलन \emph{!!{injective}} असते तेव्हा
आणि केवळ तेव्हाच, जेव्हा प्रत्येक $y \in B$ साठी~$f(x) = y$ असणारा
जास्तीत जास्त एक $x \in A$ असतो. अशा फलनाला $A$ पासून~$B$ पर्यंतचे
!!a{injection} म्हणतो.
\end{defn}

\begin{explain}
$f$ हे !!a{injection} आहे असे दाखवायचे असल्यास, $f$ च्या प्रांतातील
कोणत्याही !!{element}s $x$ आणि $y$ साठी $f(x)=f(y)$ असल्यास $x=y$
असते असे दाखवावे लागते.
\end{explain}

\begin{ex}
$f(x) = 1$ ने दिलेले स्थिर फलन $f\colon \Nat \to \Nat$ हे
!!{injective} ही नाही आणि !!{surjective} ही नाही.

$f(x) = x$ ने दिलेले एकरूपता फलन $f\colon \Nat \to \Nat$ हे
!!{injective} आणि !!{surjective} दोन्ही आहे.

$f(x) = x+1$ ने दिलेले उत्तरवर्ती फलन $f \colon \Nat \to \Nat$
हे !!{injective} आहे, पण !!{surjective} नाही.

पुढीलप्रमाणे व्याख्या केलेले फलन $f \colon \Nat \to \Nat$:
\[
  f(x) =
  \begin{cases}
    \frac{x}{2} & \text{जर $x$ सम असेल} \\
    \frac{x+1}{2} & \text{जर $x$ विषम असेल.}
  \end{cases}
\]
हे !!{surjective} आहे, पण !!{injective} नाही.
\end{ex}

\begin{explain}
अनेकदा !!{injective} आणि !!{surjective} हे दोन्ही गुणधर्म असलेल्या
फलनांचा विचार करायचा असतो. अशा फलनांना !!{bijective} म्हणतो.
ती \olref{fig:bijective} मधील फलनासारखी दिसतात. !!^{bijection}s ना
कधी कधी \emph{एकास-एक संगती} असेही म्हणतात, कारण ती सहप्रांतातील
घटकांची प्रांतातील घटकांशी एकमेव जोडी लावतात.
\begin{figure}
  \olasset{assets/diagrams/bijective.tikz}
  \caption{!!^a{bijective} फलन सहप्रांतातील घटकांची प्रांतातील
    घटकांशी एकमेव जोडी लावते.}
  \ollabel{fig:bijective}
\end{figure}
\end{explain}

\begin{defn}[!!^{bijection}]
$f \colon A \to B$ हे फलन \emph{!!{bijective}} असते तेव्हा आणि
केवळ तेव्हाच, जेव्हा ते !!{surjective} आणि !!{injective} दोन्ही असते.
अशा फलनाला $A$ पासून~$B$ पर्यंतचे (किंवा $A$ आणि~$B$ यांमधील)
!!a{bijection} म्हणतो.
\end{defn}

\end{document}
